\documentclass[11pt,letterpaper]{article}
\usepackage[margin=0.85in]{geometry}
\usepackage[T1]{fontenc}
\usepackage[utf8]{inputenc}
\usepackage{lmodern,microtype}
\usepackage{amsmath,amssymb,mathtools}
\usepackage{xcolor,booktabs,longtable,tabularx,array,enumitem}
\usepackage{listings,fancyhdr,lastpage,hyperref,xurl}
\definecolor{ink}{HTML}{16324F}
\definecolor{accent}{HTML}{207D88}
\definecolor{pale}{HTML}{F2F5F8}
\hypersetup{colorlinks=true,linkcolor=ink,urlcolor=accent,citecolor=ink,
 pdftitle={Shortening Lean proofs in ProveIt: a source-grounded refactoring study},
 pdfauthor={Prepared for Vladimir Reshetnikov}}
\setlength{\parindent}{0pt}
\setlength{\parskip}{0.5em}
\setlength{\emergencystretch}{3em}
\setlist{itemsep=3pt,topsep=4pt,leftmargin=*}
\pagestyle{fancy}
\fancyhf{}
\fancyhead[L]{\small\color{ink}ProveIt / DiophantineConsistency}
\fancyhead[R]{\small Proof simplification study}
\fancyfoot[L]{\small 26 October 2026}
\fancyfoot[R]{\small\thepage\ / \pageref{LastPage}}
\setlength{\headheight}{14pt}
\lstdefinelanguage{Lean}{morekeywords={theorem,lemma,def,by,exact,rw,simp,only,
 have,fun,intro,apply,cases,constructor,induction,with,namespace,end,import,
 variable,private,refine,simpa,using,let,forall,if,then,else},sensitive=true,
 morecomment=[l]{--},morecomment=[s]{/-}{-/},morestring=[b]"}
\lstset{language=Lean,basicstyle=\small\ttfamily,keywordstyle=\color{ink},
 commentstyle=\color{gray},breaklines=true,columns=fullflexible,keepspaces=true,
 showstringspaces=false,frame=single,rulecolor=\color{gray!30},
 backgroundcolor=\color{pale},aboveskip=7pt,belowskip=7pt}
\newcommand{\code}[1]{\nolinkurl{#1}}
\newcommand{\repo}{https://github.com/VladimirReshetnikov/ProveIt}
\newcommand{\commit}{60e9c5d620c6c387d932876cfa2d9376f64ea56c}
\newcommand{\src}[1]{\href{https://github.com/VladimirReshetnikov/ProveIt/blob/60e9c5d620c6c387d932876cfa2d9376f64ea56c/Computability/DiophantineConsistency/Lean/DiophantineConsistency/#1.lean}{\texttt{#1}}}
\newcommand{\checked}{\textbf{Checked}}
\newcommand{\proposed}{\textbf{Proposed}}
\newcommand{\CoeffStatus}{\checked: the complete patched module compiles and shrinks from 679 to 127 physical lines, a reduction of 552 lines (81.3\%). All nine public declaration headers are unchanged. Both the shared envelope theorem and the published weighted bound have axiom reports containing only \code{propext}, \code{Classical.choice}, and \code{Quot.sound}. The combined theorem uses a locally scoped \texttt{\detokenize{maxHeartbeats 2000000}} option: gathering 22 former proofs into one theorem changes resource budgeting. This is a maintained-source reduction, not a compilation-speed claim.}
\newcommand{\RotationStatus}{\checked: both exact-signature alternatives compile against the actual upstream module. Their axiom reports contain only \code{propext} and \code{Quot.sound}.}

\begin{document}
\begin{center}
{\Large\bfseries\color{ink}Shortening Lean proofs in ProveIt}\\[5pt]
{\large A source-grounded refactoring study of DiophantineConsistency}\\[7pt]
Prepared for Vladimir Reshetnikov\quad $\cdot$\quad 26 October 2026
\end{center}

\section*{Executive assessment}

The best opportunities are to share the algebraic and semantic interfaces that
many existing proofs already implement. Replacing a substantial proof family
with one invariant is more valuable than shortening individual tactic scripts.
This development already contains good examples of that design: generic DAG
node semantics, an addition-chain certificate checker, semiring gate semantics,
and a structural proof of the actual tile permutation. Several other modules
still duplicate the interfaces around those results.

This study combines an inventory of the active Lean tree, close reading of
selected proof families, compiled refactoring prototypes, and reviewable patches.
The deliverable is deliberately explicit about the difference between a complete
patched module that has been checked, a checked generic helper whose downstream
migration is still proposed, and an architectural idea that has not been implemented.
It does not establish a new consistency equivalence or complete the remaining
current-equation formalization.

Two complete checked patches are supplied: the package-52 coefficient-bound
module decreases from 679 to 127 physical lines, and the representation module
from 260 to 245. The latter saves 23 nonblank declaration lines after charging
for its shared helpers. Six additional checked modules demonstrate the reusable
interfaces and exact-signature alternatives discussed below.

The principal recommendations are:
\begin{enumerate}
\item Replace the repeated coefficient-envelope proof trees by one finite row
  split and a fixed arithmetic proof schema, preserving the sharp envelope.
\item Share the natural-assignment renaming law in representation constructors.
  The checked alternatives preserve the exact expressions and witness counts.
\item Use a scaled quotient/remainder identity for half-radix rotation, and a
  common residue API for cyclic block positions.
\item Share scanner iteration invariants and weighted finite-fiber summation.
\item Reuse the existing DAG and power-certificate interfaces to reduce generated
  proof scaffolding, measuring elaboration before removing cached declarations.
\item Reorganize expression and gate-semantics modules so each structural
  induction has one owner.
\end{enumerate}

The two generated circuit-step files contain 6,567 physical lines. This is a
large source-maintenance opportunity, but those lines are generated expansions,
not 6,567 lines of distinct mathematical argument. No whole-project percentage
reduction or compilation speedup is claimed.

\newpage
{\small\tableofcontents}
\newpage

\section{Snapshot, scope, and measurement}

All repository observations refer to commit
\href{https://github.com/VladimirReshetnikov/ProveIt/commit/60e9c5d620c6c387d932876cfa2d9376f64ea56c}{\texttt{60e9c5d620c6c38\ldots}}.
The commit was retrieved from the public repository on 26 October 2026. Its
message is ``Prove exact cell translation, coarse field offsets and cycle
exclusions.'' The project pins Lean \code{v4.32.0} and mathlib commit
\code{81a5d257c8e410db227a6665ed08f64fea08e997}; the experiments use those exact
versions. The upstream checkout was kept unchanged.

The requested topic also contains archived research and incoming material. The
inventory in this report concerns the active \code{Computability/DiophantineConsistency/Lean}
tree. Its main module directory includes some checks and audits as well as
mathematical development; it should not be described as exclusively production proofs.

\begin{center}
\begin{tabular}{lrr}
\toprule
Measure & Active Lean tree & Main module directory\\
\midrule
Lean files & 546 & 363\\
Physical lines & 73,971 & 68,308\\
Nonblank lines after comment removal & 59,832 & 55,127\\
Lexically identified theorem/lemma declarations & 4,967 & 4,902\\
Explicit \texttt{\detokenize{import Mathlib.Tactic}} lines & 8 & 8\\
\code{maxHeartbeats} option commands & 134 & 134\\
\bottomrule
\end{tabular}
\end{center}

These are reproducible lexical measurements, not counts of independent
mathematical results or kernel proof-term sizes. The script removes nested Lean
comments and preserves line boundaries. Declaration spans are estimates, because
the script is not a Lean parser. Source hashes and per-file data are in
\code{inventory/files.csv}; the script is \code{scripts/audit_inventory.py}.

The lexical scan found no \code{sorry}/\code{admit} tokens, explicit \code{axiom}
commands, or \code{native_decide} tokens in the active tree. This is narrower
than a transitive axiom audit. The supplied new proofs have their own actual
\texttt{\#print axioms} outputs; the entire existing development was not rebuilt
or transitively audited for this report.

\subsection{What counts as a useful simplification?}

Four costs should be tracked separately: maintained source, mathematical proof
structure, elaboration cost, and the transparency of the resulting theorem.
The practical objective is a smaller set of reusable invariants with explicit
interfaces. A one-line tactic that unfolds an enormous term can have worse
elaboration behavior and be harder to repair than a longer sequence of named lemmas.

Exact expression syntax and witness counts matter here. A merely equivalent
Diophantine representation can change the constants, size bounds, executable
recipe, or later calibration. The representation refactor therefore includes
equalities between the alternative constructors and the original constructors,
proved by reflexivity. The coefficient refactor retains the literal coefficient
function and the original envelope. The arithmetic interfaces keep all essential
positivity and cyclic-domain assumptions.

\section{Checked refactor: assignment naturality}

\subsection{The repeated argument}

\src{Representation} repeatedly establishes that renaming input and witness
indices commutes with the natural-to-integer assignment. The same calculation
appears in \code{Rep.reindex}, twice in \code{Rep.conjunction}, and twice in
\code{Rep.disjunction}. The underlying law is
\[
  \operatorname{natAssignment}(v,w)\circ\operatorname{Sum.map}(f,g)
  =\operatorname{natAssignment}(v\circ f,w\circ g).
\]
Its proof is function extensionality followed by the two constructors of a sum.
Composing it with the existing \code{Expr.eval_rename} produces the semantic
interface used at every call site.

\begin{lstlisting}
theorem natAssignment_comp_sumMap (v : tau -> Nat) (w : beta -> Nat)
    (f : sigma -> tau) (g : alpha -> beta) :
    Function.comp (natAssignment v w) (Sum.map f g) =
      natAssignment (Function.comp v f) (Function.comp w g) := by
  funext i
  cases i <;> rfl
\end{lstlisting}
The complete, compilable source is \code{Lean/ProofSimplification/RepresentationNaturality.lean}.

The reindex correctness proof becomes an explicit simplification using the
shared evaluation lemma. Conjunction and disjunction retain their existing
\code{Fin.addCases} witness assembly. Their mathematical constructions are
unchanged; only the repeated assignment calculation is centralized.

\subsection{Measured result and compatibility}

\begin{center}
\begin{tabular}{lrr}
\toprule
Declaration & Original nonblank lines & Refactored\\
\midrule
\code{reindex} & 12 & 5\\
\code{conjunction} & 29 & 16\\
\code{disjunction} & 30 & 17\\
Shared helper declarations & 0 & 10\\
\midrule
Total & 71 & 48\\
\bottomrule
\end{tabular}
\end{center}

The net saving in those declarations is 23 nonblank lines, or 32.4\%.
Including comments, spacing, and helper namespace scaffolding, the complete
patched module has 245 physical lines instead of 260. Both measurements are
reported so that omitted overhead is not mistaken for a saving.

\checked: the additive prototype, the original complete module, and the
complete patched module compile. The alternative \code{reindex},
\code{conjunction}, and \code{disjunction} constructors equal the original
constructors by \code{rfl}. Their stored syntax and witness count are
definitionally identical; proof irrelevance identifies the correctness fields.
The patch preserves public names. \texttt{\detokenize{git apply --check}} succeeds against
the pinned checkout.

The shared lemma is invoked explicitly. It is not necessary to make large
expression definitions global simplification rules. The patch and the complete
patched file are included separately, so reviewers can inspect either the
change or the final module.

\section{Coefficient envelopes: one proof schema for all rows}

\subsection{Why this is a substantial local target}

\src{NativeCoefficientBounds} has 679 physical lines. Its documented subject
is the package-52 native parent. Thus the largest measured local reduction
belongs to that existing arithmetic layer; the current length-2011 calibration
has separate field and mask proofs. A large part of the file
is the family \code{rowEnvelope_0} through \code{rowEnvelope_21}, followed by a
public wrapper. Each private proof repeats three positivity facts and follows
the conditional coefficient recipe through a tree of absolute-value bounds.
The repeated region contains 158 occurrences of the same restricted
\texttt{\detokenize{linarith only}} closure and 136 applications of \code{abs_ite_le}.

Write the four ports as $b,t,m,f$, where $t$ is the large transport parameter.
The exact envelope is
\[
 E=2b^2+2fb+m+f+2b+tb+t+8.
\]
Given $b,t,m,f\geq0$, the proof uses only those four inequalities and
$b^2\geq0$, $fb\geq0$, $tb\geq0$. Once a coefficient branch is exposed, both
sides of its absolute-value inequality are linear combinations of these
quantities. Restricted linear arithmetic is therefore sufficient.

\subsection{The proposed replacement structure}

Keep the original definition of \code{envelope} and the actual
\code{rowCoefficient}. Introduce the three shared nonnegativity facts once.
Split the 22-valued row index, reduce that row's definition, and traverse its
conditional tree with the existing semantic helper \code{abs_ite_le}.
One \texttt{\detokenize{repeat' apply abs_ite_le}} replaces the manually written branch tree.
At every leaf, rewrite \code{abs_le} and solve the two inequalities with the
same fixed collection of facts.

This removes proof duplication while retaining the sharp dependence on the
transport port. In particular, the refactor does not replace the envelope by a
convenient but larger quadratic expression in $t$. The downstream
\code{envelope_lt} estimate and correction bound keep their mathematical content.

\CoeffStatus

The shared proof body is:
\begin{lstlisting}
have hsq := sq_nonneg (ports 0)
have hmf := mul_nonneg h3 h0
have hkp := mul_nonneg h1 h0
fin_cases k <;>
  simp only [rowCoefficient, envelope,
    Matrix.cons_val_zero', Matrix.cons_val_succ'] <;>
  repeat' apply abs_ite_le
all_goals
  rw [abs_le]
  constructor <;>
    linarith only [h0, h1, h2, h3, hsq, hmf, hkp]
\end{lstlisting}

The supplied patch has a complete module copy, not just an isolated proof
fragment. Its public theorem has the same quantifiers, hypotheses, and conclusion.
The code-generating origin of coefficient data does not make the arithmetic
schema an oracle: branch reduction and \code{linarith} produce terms checked
by Lean's kernel.

\subsection{Maintenance considerations}

The choice of branch interface matters: a trial using generic
\code{split_ifs} reached a simplification-step limit on the longest row.
Applying the shared absolute-value lemma handles both branch values without
evaluating the pair-index equality tests or carrying their branch hypotheses.
This is a concrete example where a small semantic lemma makes the concise
proof more robust than unrestricted case splitting.

This proof depends deliberately on the finite recipe representation. If a row
changes, Lean rechecks its branches against the same envelope and fails if the
old bound no longer holds. Use \texttt{\detokenize{linarith only}} with the fixed sign facts,
so branch equalities do not become accidental proof dependencies. A named local
tactic can be helpful if several coefficient families use the identical
schema, but adding a general tactic framework for this one family is unnecessary.

A further, unimplemented improvement is to represent the repeated coefficient
values by a small finite catalog. The baseline's 136 nondefault branches have
25 distinct textual value expressions, plus the default zero. The row/pair
lookup would return a catalog
entry, and its envelope bound would be proved once per entry. A separate exact
evaluation theorem must identify the catalog lookup with the existing literal
coefficient function. That extra equality is what makes the catalog a proof
refactor rather than a changed arithmetic encoding. This is a larger generator
and data-layout change than the immediate proof-schema patch.

\section{Scaled division and cyclic residue interfaces}

\subsection{An exact identity removes finite digit enumeration}

For natural numbers $s,x,y$ with $s>0$,
\begin{equation}\label{eq:scaled}
 (sx)\bmod(2s)+\left\lfloor\frac{sy}{2s}\right\rfloor
 =s(x\bmod2)+\left\lfloor\frac y2\right\rfloor.
\end{equation}
This follows directly from \code{Nat.mul_mod_mul_left} and
\code{Nat.mul_div_mul_left}. It is an equality for arbitrary $x,y$; no digit
bound is required. The assumption $s>0$ is essential: at $s=0$ the left side
is zero, whereas the right side need not be.

If $y\leq3$, then $x\bmod2\leq1$ and $y/2\leq1$, giving
\[
 (sx)\bmod(2s)+(sy)/(2s)\leq s+1.
\]
If $s$ is even as well, the parity of the same expression is exactly $y/2$.
These are \code{scaled_half_bound} and \code{scaled_half_parity} in the supplied
\code{ArithmeticRefactors} module. Notice that the bound on $x$ is unnecessary;
only the predecessor value $y$ needs its small-digit bound.

\src{NativeCyclicRotation} currently proves \code{shiftDigits_half_bound}
using four cases for the current digit and a separate quotient estimate.
Its \code{shiftDigits_half_parity_89} uses the sixteen combinations of two
four-valued digits. Equation~\eqref{eq:scaled} explains both arguments at once.
The supplied \code{RotationRefactors} alternatives use the same original theorem
signatures and retain the absence of a \texttt{\detokenize{NeZero n}} assumption. They do not
call the old bound or parity theorem.

\RotationStatus

The parity proof is not necessarily shorter as an isolated four-line script:
the original enumeration is already compactly written. The improvement is an
explicit reusable reason for the result, with less dependence on the particular
radix exponent or digit case count. Shared-helper overhead should be included
when assessing a one-off replacement.

\subsection{Factor the correct cyclic residue statements}

Near-identical block-residue proofs appear in \src{CurrentClause2011},
\src{CurrentReflectedGeometry2011}, \src{NativeLocalMask54}, and
\src{NativeCleanBands54}. The useful common parameter is the block length $L$,
with positions in $\mathbb Z/(NL)\mathbb Z$. The checked generic identities are
\[
 \operatorname{val}(b+e)\bmod L=(\operatorname{val}(b)+e)\bmod L,
\]
\[
 (\operatorname{val}(k-e)+e)\bmod L=\operatorname{val}(k)\bmod L.
\]
Both follow by casting into the full cyclic modulus, then restricting modular
equality along $L\mid NL$. The supplied statements use \texttt{\detokenize{[NeZero (N * L)]}}.

The subtraction identity is intentionally written with addition of $e$ after
taking the representative. Replacing cyclic subtraction by truncated natural
subtraction would break wrapped positions. A common arithmetic module can
export these facts below both the current and legacy calibration layers.
\src{CurrentCellTranslation2011} can then consume them while retaining its
explicit wrapped and unwrapped source-position theorems.

\checked: the generic residue and scaled-arithmetic lemmas compile. Migrating
every current and legacy caller is still a proposed follow-up. Their common
arithmetic does not identify the different masks, offsets, or tile calibrations.

\section{Shared scanner machinery without changing raw-input soundness}

\src{AnnotatedAxiomContext}, \src{AnnotatedBoundScan},
\src{AnnotatedLineScan}, \src{AnnotatedContext}, and
\src{AnnotatedRenameScan} all define a numeric scanner by primitive recursion
on fuel. The transition function may have additional fixed parameters, but
after those are supplied it is an endofunction on states.

The checked \code{ScannerCore} defines exactly the existing recursion pattern:
\begin{lstlisting}
def scan (step : S -> S) (fuel : Nat) (initial : S) : S :=
  fuel.rec initial (fun _ acc => step acc)
\end{lstlisting}
It proves the successor law, preservation of a fixed state, preservation of an
invariant, and preservation of ``reject or good.'' All four axiom reports are
empty. Existing scanner definitions can remain unchanged and use these facts
by definitional equality; a supplied example checks the raw recursion directly.

For the soundness interface, let $r$ be the rejecting state and suppose
\[
 T(r)=r,\qquad G(s)\Longrightarrow T(s)=r\ \lor\ G(T(s)).
\]
Then starting from a good state, every fuel-bounded iterate is rejecting or
good. This is the invariant $P(s):=(s=r\lor G(s))$ under the ordinary iteration
theorem. The disjunction need not be exclusive.

For a concrete proof checker, $G$ can say that a raw numeric state encodes an
actual history of provable judgments. The transition proof must still establish
that invariant from the local checker result. It must not assume that arbitrary
raw inputs already parse as valid formulas or derivations. Likewise, the helper
does not prove termination, enough fuel, successful acceptance, or the
primitive-recursive realization of the scanner program. Those obligations
remain in the concrete modules.

The first migration should target \code{scan_zero}, \code{scan_succ}, and
\code{scan_done}; these are low-risk direct uses. Only then should the common
shape of \code{scan_sound} be factored, with the raw-input theorem and explicit
program code checked again. A change to a program constructor requires separate
attention to the existing code-size claims.

\section{One finite-fiber identity for folded stencils}

\src{NativeCyclicConvolution} proves \code{mass_foldStencil},
\code{convolution_foldStencil}, and the finite-sum step inside
\code{word_foldStencil} by repeatedly swapping sums and collapsing an equality
filter. All are instances of
\begin{equation}\label{eq:pushforward}
 \sum_{k\in K}\left(\sum_{i\in I:p(i)=k}a_i\right)w_k
 =\sum_{i\in I}a_iw_{p(i)}.
\end{equation}
The supplied \code{FinitePushforward.sum_fiber_mul} proves this for finite
index types, an arbitrary map $p:I\to K$, and a semiring. Multiplication need
not be commutative, because the order $a_iw_{p(i)}$ is retained. Decidable
equality on $K$ is used to express the finite filter.

Three specializations are immediate: $w_k=1$ for mass,
$w_k=c(q-k)$ for convolution at $q$, and $w_k=\rho^{\operatorname{val}(k)}$
for a word. For the last specialization, the existing
\code{pow_mod_length} argument is still required to replace the reduced
exponent by the original exponent modulo $\rho^n-1$.

\checked: the generic identity compiles. It assumes no injectivity of $p$,
so all collisions are retained. It does not establish no-carry bounds, mask
compatibility, or uniqueness of a word representation. A safe integration
replaces the repeated finite-sum rearrangement and preserves those separate
arithmetic obligations.

This is a useful example of a small abstraction with several interpretations.
It makes the collision behavior obvious at the interface instead of requiring
each proof reader to reconstruct it from a double sum.

\section{Generated circuits: reduce proof scaffolding with existing checkers}

\subsection{Replace repetitive node interfaces selectively}

\begin{center}
\begin{tabular}{lrrr}
\toprule
Module & Physical lines & Node equations & Power equations\\
\midrule
\code{GapCircuitSteps66} & 3,122 & 484 & 86\\
\code{SplitCircuitSteps59} & 3,445 & 515 & 107\\
\midrule
Total & 6,567 & 999 & 193\\
\bottomrule
\end{tabular}
\end{center}

Both files explicitly identify themselves as generated. Their local node
equations already apply \code{ArithmeticDAG.run_node_equation} from
\src{ArithmeticDAGSemantics}. Thus the mathematical invariant is already shared;
the remaining duplication is the interface and its generated normalization.

\proposed: expose one theorem indexed by
\texttt{\detokenize{Fin naturalNodes.length}}, giving the local equation for that node.
At a use site, instantiate it and reduce only the selected node lookup.
Alternatively, a small proof-producing elaborator can perform those two steps.
The generator can retain named theorems at semantic boundaries such as the
base, a source quantity, or an output residual, while avoiding an independent
declaration for every routine node.

There is a concrete tradeoff. Existing named node equations cache their reduced
types. Replacing them by repeated lookup reduction against a large literal list
can make elaboration slower even if the source is shorter. A measured migration
should compare the downstream recipe modules, not just the step-file line count.
Compatibility aliases and newly emitted data must be counted in the final size.

\subsection{The power-label checker is already available}

\src{ArithmeticDAGPowers} contains \code{PowerStep}, \code{run_power}, and
\code{value_eq_of_powerRewrite}. These establish soundness of a table of
exponent labels for the actual topological circuit. The generation scripts
already compute labels of this form: the base receives label one, and a
multiplication of two labelled predecessors receives the sum of their labels.

\proposed: emit the label table as data, prove its finite local obligations by
kernel-checked computation, and apply \code{run_power} once. Named powers can
then be short specializations. This is preferable to inventing a second
symbolic-power correctness mechanism or repeatedly invoking \code{ring} after
each predecessor rewrite.

The 193 generated power equations identify the candidate population; they are
not a promised 193-theorem net saving. The label data and API aliases have a
cost. The proof must retain the actual circuit, its topological order, the
correct base, operand labels, and index bounds. This migration has not been
implemented in the supplied patches.

\section{Expression and compiler architecture}

\subsection{Centralize residual-square semantics and degree bounds}

\src{DAGQuartic}, \src{NaturalDAGQuartic}, and
\src{QuadraticRegisterSystem} repeat an evaluator lemma for
\code{QuarticSyntax.sumExpressions} and a theorem that a finite sum of squared
integer residuals vanishes exactly when every residual vanishes. Their proofs
already use the library's \code{Finset.sum_mul_self_eq_zero_iff}.

\proposed: put list evaluation and the finite residual-square interface below
the three compiler layers. Each compiler then supplies only its own residual
family. A companion degree theorem should derive a quartic upper bound from
quadratic bounds on the residuals. This avoids reproving both the semantic
and polynomial-size bookkeeping for every witness-layout variant.

The zero-set theorem needs its mathematical domain. Over $\mathbb Z$ the
nonnegative squares cannot cancel. Over an arbitrary commutative ring the
statement can fail, so a ring-generic evaluator theorem is not by itself a
ring-generic residual-zero theorem. Use the existing ordered/integer interface
or state exactly the additional assumptions that a broader theorem needs.

\subsection{Move basic syntax facts below the consumers}

\code{QuarticCompiler.Gates.polynomial_rename} and
\code{NaturalRegisterRelation.rename_toPolynomial} have the same statement:
\[
 (\operatorname{rename}(f,e)).\operatorname{toPolynomial}
 =\operatorname{MvPolynomial.rename}(f,e.\operatorname{toPolynomial}).
\]
Both are structural induction proofs. A canonical theorem in the \code{Expr}
layer removes duplicate ownership; aliases can retain existing call sites.

Similarly, \code{QuadraticRegisterSystem.syntaxDegree} and
\code{polynomial_degree_le} concern expression syntax, not register execution.
Moving them into a low-level expression module allows all residual compilers
to use one syntactic degree bound. Such a bound is an upper bound; it must not
replace a semantic identity or an exact degree claim that depends on cancellation.

\subsection{Share the gate soundness and completeness inductions}

\src{QuarticCompiler} contains natural/integer gate soundness and completeness
proofs. \src{GatesSemiring} contains parallel generic semiring proofs for the
same wire structure, root, and residual syntax. The existing import direction
is important: \code{GatesSemiring} imports \code{QuarticCompiler}.

\proposed migration:
\begin{enumerate}
\item Extract \code{Wire}, \code{root}, \code{count}, and \code{residualSyntax}
  into a dependency-light \code{GatesSyntax} module.
\item Put generic semiring evaluation, soundness, and completeness above that syntax.
\item Derive the natural/integer versions using existing
  \code{NatExpr.evalSemiring_nat} and
  \code{Expr.eval_eq_zero_iff_splitNatural}.
\item Retain the executable constructors and verify exact witness and operation counts.
\end{enumerate}
Simply calling a generic gate theorem from the current compiler without first
extracting syntax would create an import cycle. The real simplification is one
owner for each induction and short transport results at the boundaries.

\subsection{Respect the object-language boundary}

The \code{Object*} modules prove arithmetic facts inside the repository's
sealed ZFC calculus. \src{ObjectNaturalSemiring} ranges over the model's entire
internal natural numbers, including any nonstandard elements.
\src{ObjectResidualSquares} uses this interpretation to obtain an internal
residual-conjunction statement.

A proof about ordinary Lean naturals alone cannot replace such a theorem.
The useful refactor is to identify a generic algebraic statement that holds
for the already constructed internal semiring, then use the explicit semantic
and syntactic transport theorems. Avoid replacing internal assignments by
standard numerals or treating an external semantic proof as an internal
derivation. This boundary is central to preserving the intended consistency claim.

\section{Further targeted opportunities}

\subsection{Bundle carry witnesses with named fields}

\code{NativeWindowCarry.source_convolution_carries} and
\code{CurrentWindowCarry2011.source_carries} return carry and digit functions
with several nested facts: carry bounds, digit bounds, quotient recurrence,
the carrying equation, the reconstructed word, and mask compatibility.
Repeatedly unpacking these long conjunctions increases argument-management work.

\proposed: define a small record for the generic carry solution, with named
fields for the recurrence and bounded carries, then put source-specific word
and mask facts in a separate wrapper if needed. Where a digit is definitionally
the remainder, its bound and carrying equation should be derived once from
division with a positive radix. A record improves theorem navigation and makes
it harder to supply a fact about the wrong digit array.

The record must not collapse local carry reset into a global no-carry claim.
Nor should it introduce a hidden temporal-alignment premise. Those distinctions
are present in the current arithmetic and remain essential. This proposal is
about a clearer interface; no line reduction is measured for it yet.

\subsection{Use standard membership lemmas for small inductions}

The private \code{AnnotatedProofChecker.mem_of_lookup} repeats a list/index
induction to infer $j\in\mathrm{seen}$ from
\texttt{\detokenize{seen[i]? = some j}}. The pinned core library already provides
\code{List.mem_of_getElem?}. The supplied \code{ListLookup} file checks the
corresponding generic theorem directly, with only \code{propext} in its axiom report.

This is a checked API-level cleanup, not a replayed patch to the entire annotated
checker. Its surrounding natural-deduction rule proof has genuinely distinct
cases. Explicit rule constructors can remain visible while standard list facts
discharge the repeated premise-membership obligations.

\subsection{Use existing structural proofs as patterns}

Two recent modules already illustrate the intended direction.
\src{WangTiles942Permutation} proves a structural left inverse for the actual
tile table, then proves injectivity and rank properties. It does not perform a
quadratic comparison of every tile pair. \src{NativeFieldCycles2011} uses
periodicity and modular arithmetic to exclude the remaining offsets, including
a short two-step residue obstruction where no finiteness is needed.

These are existing achievements, not new results of this audit. They suggest a
method for future proofs: identify the structural inverse, invariant, or local
residue obstruction before generating a large finite case analysis. In
\src{CurrentKernel2011}, small batches of kernel-computed finite sums also serve
as deliberate elaboration boundaries. Removing those boundaries to obtain one
large decision call needs performance evidence.

\section{Implementation order and validation policy}

\begin{longtable}{>{\raggedright\arraybackslash}p{0.06\textwidth}>{\raggedright\arraybackslash}p{0.39\textwidth}>{\raggedright\arraybackslash}p{0.20\textwidth}>{\raggedright\arraybackslash}p{0.24\textwidth}}
\toprule
Stage & Change & Evidence in this package & Next integration step\\
\midrule\endhead
1 & Assignment naturality & Complete checked patch & Apply in a branch; replay direct consumers\\
1 & Coefficient envelope schema & Complete checked patch & Check downstream coefficient-bound consumers\\
1 & Residue and scaled arithmetic & Checked generic helpers & Migrate callers with original signatures\\
1 & List membership alias & Checked library API & Replay annotated checker\\
2 & Scanner fixed-point/iteration facts & Checked generic helper & Replace basic lemmas before soundness invariants\\
2 & Weighted finite fibers & Checked generic helper & Specialize mass, convolution, word sum step\\
2 & Expression renaming and residual-square API & Source-grounded proposal & Move low-level lemmas; retain aliases\\
3 & Gate semantics shared over semirings & Architecture proposal & Extract syntax first; avoid import cycle\\
3 & Indexed circuit steps and power labels & Existing checker identified & Pilot one recipe; measure elaboration and memory\\
3 & Carry record and generic block geometry & Interface proposal & Keep explicit source and phase obligations\\
\bottomrule
\end{longtable}

For a proof-only patch, review the theorem statement and executable data first.
Then compile the complete changed module against the pinned dependencies,
inspect its axiom report, and replay the direct consumers that depend on its
API or reduction behavior. For a structural change, add an exact equality or
transport theorem when data preservation is part of the claim. Source-level
shortening by itself is not evidence that a witness representation is unchanged.

The included experiments check complete patched modules independently, so a
candidate is not accidentally accepted by importing the old theorem with the
same name. Additive prototypes can import old modules for definitions, but their
proofs must not invoke the old theorem they purport to replace. The rotation
alternatives meet this condition.

\subsection{Trusted computation and simplification boundaries}

The existing development uses kernel-checked finite computation. The suggested
certificate approach retains that property. Lean's documentation distinguishes
\texttt{\detokenize{decide +kernel}} from \code{native_decide}/\texttt{\detokenize{decide +native}}, which
uses native evaluation and changes the trusted dependencies of a proof
\cite{lean-tactics,lean-validate}. This report does not propose changing that
trust boundary to save source lines.

Use explicit semantic rewrite lemmas around representations, scopes, and
assignments. Lean's \texttt{\detokenize{simp only}} permits a deliberately bounded rewrite set
\cite{lean-simp}; it is useful at those interfaces. Broad simplification is
often appropriate for small constructor goals, but expanding an entire circuit
or coefficient recipe as a global simp rule can obscure the intended boundaries.

\section{Validation results and reproducibility}

The bundle contains source hashes, patch checks, compilation logs, and
machine-readable records of successful runs. The generic residue, scaled
arithmetic, finite-fiber, scanner, assignment, and list-lookup interfaces are
compiled with the pinned toolchain. The detailed status of each final source
is recorded in \code{validation/final-validation.json}.

The reported axiom sets for the new algebraic and representation proofs are
subsets of \{\code{propext}, \code{Classical.choice}, \code{Quot.sound}\}.
The four scanner lemmas require no axioms. No new \code{sorry}, custom axiom,
or native-evaluation oracle is used. A separate source review checked the
packaged generic interfaces for hidden hypothesis strengthening and circularity;
that review is included as an evidence note and is not described as a second
compilation run.

The complete DiophantineConsistency library and all downstream consumers were
not rebuilt. The generic scanner and finite-sum abstractions have not been
installed throughout the project. Architectural proposals remain proposals.
These limits are reflected in the status labels rather than being hidden by
the existence of a successful small example.

Individual compilation times are retained for reproducibility, not interpreted
as performance benchmarks. For example, one representation replay took 3.90 s
for the original and 5.01 s for the patched file under concurrent load. Those
observations do not demonstrate a speedup. Other checks ran under substantial
shared CPU and disk pressure. A serious performance comparison should isolate
the same complete target, use the same cached imports and output mode, repeat
runs, and report peak memory as well as elapsed time.

\subsection{How to use the bundle}

Start from the pinned ProveIt checkout with its Lean dependencies available.
The \code{README.md} describes the exact commands. Run
\code{scripts/verify.py} against that checkout to check patch applicability,
compile the prototypes and full patched copies, and record a fresh validation
result. The script leaves the checkout's source files unchanged. It expects
the listed baseline dependency modules to have already been built; this avoids
silently launching a full project build.

The original source remains authoritative. The two unified patches are intended
for review on an integration branch. The \code{patched/} copies make the
complete proposed modules easy to inspect. The full local toolchain, package
caches, and compiled objects are deliberately not included in the archive;
the pinned versions and commands are sufficient to recreate them.

\section{Further research questions}

\begin{enumerate}
\item Can one small certified arithmetic-recipe interface cover the repeated
  coefficient families while retaining their distinct sharp bounds? The useful
  abstraction would expose a finite set of coefficient expressions and an
  explicit sign basis, rather than treating the final inequality as an opaque tactic call.
\item What is the best granularity of cached circuit equations? An indexed
  theorem minimizes declarations, while named intermediate equations cache
  normalization. A hybrid selected by actual downstream use may be better than
  either extreme.
\item Can the existing exponent-label checker be extended to additional
  arithmetic invariants, such as syntactic degree or coefficient support, with
  one shared local soundness theorem for circuit folds?
\item How much of the annotated checker can be expressed through one scanner
  invariant without obscuring arbitrary-raw-input soundness or changing the
  primitive-recursive code whose size is counted?
\item Which internal-natural algebra results can be transported through the
  established semiring interpretation, and which depend essentially on the
  object-language encoding and sealing machinery? A precise inventory would
  prevent both duplicate inductions and invalid external-to-internal shortcuts.
\item Can a parameterized cyclic-block library serve both the legacy and current
  arithmetic without making calibration-dependent assumptions implicit? The
  shared residue results are a small verified first step; masks and field
  positions should remain explicit data with proved interfaces.
\end{enumerate}

\begin{thebibliography}{9}
\bibitem{proveit}
V. Reshetnikov et al., \emph{ProveIt, DiophantineConsistency}, pinned commit
\href{https://github.com/VladimirReshetnikov/ProveIt/tree/60e9c5d620c6c387d932876cfa2d9376f64ea56c/Computability/DiophantineConsistency}{60e9c5d620c6c38\ldots}.
Module links throughout this report refer to this snapshot.
\bibitem{toolchain}
ProveIt, \href{https://github.com/VladimirReshetnikov/ProveIt/blob/60e9c5d620c6c387d932876cfa2d9376f64ea56c/lean-toolchain}{lean-toolchain}
and \href{https://github.com/VladimirReshetnikov/ProveIt/blob/60e9c5d620c6c387d932876cfa2d9376f64ea56c/lake-manifest.json}{lake-manifest.json}.
\bibitem{mathlib}
The mathlib community, \href{https://github.com/leanprover-community/mathlib4/tree/81a5d257c8e410db227a6665ed08f64fea08e997}{mathlib at 81a5d257c8e410db\ldots}.
\bibitem{lean-tactics}
Lean Language Reference, \href{https://lean-lang.org/doc/reference/latest/Tactic-Proofs/Tactic-Reference/}{Tactic Reference}.
Retrieved 26 October 2026. Version-specific claims in the experiments are
additionally checked by the pinned compiler.
\bibitem{lean-validate}
Lean Language Reference, \href{https://lean-lang.org/doc/reference/latest/ValidatingProofs/}{Validating a Lean Proof}.
Retrieved 26 October 2026.
\bibitem{lean-simp}
Lean Language Reference, \href{https://lean-lang.org/doc/reference/latest/The-Simplifier/Invoking-the-Simplifier/}{Invoking the Simplifier}.
Retrieved 26 October 2026.
\end{thebibliography}

\end{document}
